% =============================================================================
%  3. Quantitative analysis  --  CORRECTED against code/global/ (2026-08-22)
%  Notation follows 02-model.tex.  Two deliberate notation changes are flagged
%  in the accompanying change list:
%    (i)  the intermediation wedge is written \bar\varsigma, not \bar\sigma,
%         because \sigma_X is risk aversion and \sigma_s the risk-process
%         innovation s.d. in 02-model.tex;
%    (ii) the entrant transfer is \omega_X throughout (02-model.tex writes
%         \omega_X^f in eq. (nw-lom) and \omega_X in eq. (bank-div)).
%  Bibliography keys used: bocola2016pass, gertler2011model, gertlerkiyotaki2010,
%  Hatchondo2009, neumeyer2005business, jermann1998asset, bohn1998behavior,
%  zettelmeyer2013, rouwenhorst1995.
%  [TBC] marks a number whose source must be filled before circulation.
% =============================================================================
\section{Empirical Analysis}
\label{sec:analysis}

The model is calibrated at a quarterly frequency to the euro-area
periphery--core pair on the eve of the sovereign debt crisis, with country $D$
standing for Greece and $F$ for the core. Three groups of parameters should be
kept apart. Most are set directly, from the literature or from pre-crisis
euro-area data. The agency friction of the intermediary block is not assigned but
\emph{solved}, so that the steady state reproduces a leverage ratio and an
intermediation wedge. The remainder are normalisations or equilibrium outcomes,
the household discount factor above all, which is solved so that private saving
clears the deposit market at the calibrated risk-free rate.

\subsection{Calibration strategy}
\label{sec:calibration}

Throughout, the two countries carry \emph{identical} parameters; asymmetry
enters through shocks alone, in that $D$'s sovereign is default-risky and $F$'s
is safe, and only $D$ is shocked. Symmetry is imposed rather than incidental: it
delivers $\bar p=1$ with zero cross-border deposit and net foreign asset
positions, and since $p$ is only weakly identified by external balance at a trade
elasticity below one, an asymmetric steady state moves $\bar p$ off unity and
opens a residual in the goods market. Two normalisations complete the picture:
$\chi_X$ is set so that $\bar N_X=1$ and $\bar Z_X$ so that $\bar Y_X=1$, so all
quantities below are shares of quarterly steady-state output. Parameter values
and empirical targets are collected in the appendix.

\paragraph{Households.}
Four features of \eqref{eq:hh-objective}--\eqref{eq:hh-borrowing} map the
functional forms directly into calibration choices. First, GHH preferences
eliminate wealth effects on labour supply: the intratemporal condition is
$\chi_X n_{it}^{1/\nu_X}=(w_{X,t}/P^{c}_{X,t})e_{it}$, so $\nu_X$ is the Frisch
elasticity and $\chi_X$ a pure scale normalisation. We set $\nu_X=2$, the value
implied by the inverse Frisch elasticity of $0.5$ in \citet{bocola2016pass}, and
normalise $\chi_X$ to $\bar N_X=1$. Second, we set $\sigma_X=1$, so that
period utility is logarithmic in the consumption--labour composite $x_{it}$;
this is the preference specification of \citet{bocola2016pass}, and it is the
convention under which the bankers' kernel discussed below satisfies
$\beta^{\rm int}(1+\bar r)=1$ at the steady state. Third, idiosyncratic labour productivity follows an
AR(1) in logs with persistence $\rho_e=0.9$ and innovation standard deviation
$\sigma_e=0.2$, discretised by the \citet{rouwenhorst1995} procedure into an
$n_e$-state Markov chain $(\mathsf{e},\Pi_X)$ whose grid is normalised so that
$\sum_e\pi^{\ast}(e)\,\mathsf{e}(e)=1$, where $\pi^{\ast}$ is the stationary
distribution of $\Pi_X$. We use $n_e=2$: what the household block has to deliver
here is an aggregate deposit supply schedule, not the shape of the wealth
distribution, and a coarse chain keeps it cheap inside a global projection solve.
Fourth,
with $\underline a_X=0$ the stationary wealth distribution is pinned down jointly
by $\beta_X$ and $\{e_{it}\}$; assets are discretised on a $250$-point grid,
curved towards the borrowing constraint.

The discount factor $\beta_X$ is \emph{not} calibrated to a liquid-wealth target.
It is solved, in the second stage of the steady state, so that aggregate
household savings equal the intermediaries' deposit funding need at the
calibrated deposit rate,
\begin{equation}
A_X(\beta_X;\bar r)\;=\;\overline{\mathrm{dep}}_X\;=\;(\bar\theta_X-1)\,\bar n_X ,
\label{eq:beta-clearing}
\end{equation}
with the right-hand side delivered by the intermediary block. This gives
$\beta_D=\beta_F=0.9945$, hence $\beta_X(1+\bar r)=0.9975<1$ as incomplete markets
require. The implied liquid-wealth ratio is therefore an output, not a target,
and it is large --- deposits are $2.07$ times annual output --- because deposits
are the only store of value and intermediaries hold the entire capital stock. It
is not comparable to a household-survey measure of liquid wealth; matching such a
target instead would make the deposit rate the equilibrium object.

\paragraph{Production and capital.}
The production block takes standard values: a capital share $\alpha_X=0.30$, a
quarterly depreciation rate $\delta_X=0.025$, and a retail demand elasticity
$\epsilon_X=6$, so that the flexible-price markup is $20\%$ and real marginal
cost is $mc_X=5/6$. The one non-standard parameter is the elasticity of Tobin's
$q$ with respect to the investment rate, $\xi_X=0.42$, which we take from the
posterior mean of \citet{bocola2016pass} rather than from a calibration target of
our own; it governs how much of a shift in the intermediaries' required return on
capital shows up as a price rather than as a quantity.

The two remaining adjustment-cost coefficients, $\gamma_{0,X}$ and
$\gamma_{1,X}$, are not free. We discipline them by requiring that adjustment
costs be \emph{inactive} at the non-stochastic steady state: the price of
installed capital is exactly one, $\bar Q_X=1$, and steady-state investment
exactly replaces depreciation, $\bar\iota_X=\delta_X$, so that
$K_{X,t+1}=K_{X,t}$. Imposing both conditions on \eqref{eq:capital-lom} and
\eqref{eq:tobins-q} pins the coefficients as functions of $\delta_X$ and $\xi_X$
alone:
\begin{equation}
\gamma_{0,X}=\frac{\delta_X^{\;\xi_X}}{1-\xi_X},
\qquad
\gamma_{1,X}=-\frac{\delta_X\,\xi_X}{1-\xi_X}.
\label{eq:jermann-coeffs}
\end{equation}
The steady state that results is a check on the block rather than a target of it:
at $\bar r^{k}=0.32\%$ per quarter (see below) the implied capital--output ratio
is $2.22$ in annual terms and the investment share is
$\delta_X\bar K_X/\bar Y_X=22.2\%$, against a Greek average of $24.3\%$ of GDP
over $1999$--$2007$.\footnote{FRED series \texttt{GRCGFCFQDSMEI} and
\texttt{GRCGDPNQDSMEI}, both in current prices.}

Firms pre-finance a fraction $\zeta_X$ of the wage bill, \eqref{eq:loan-demand};
we set $\zeta_X=1$. Since \eqref{eq:labour-demand} is the only channel from the
credit spread to output on impact, $\zeta_X$ scales the impact output response
one-for-one and $\zeta_X=0$ nests the model without it exactly. It leaves bond
prices, net worth and the spread untouched on impact, so the model's robust
predictions are not the ones this parameter governs.

\paragraph{Trade.}
The Armington aggregator \eqref{eq:ces-aggregator} carries a home bias
$\varpi=0.85$ and a trade elasticity $\eta=0.5$. The import share of $15\%$ is in
line with the euro-area periphery, and the trade elasticity is in the range
conventional in the international real business cycle literature, which is well
below the elasticities estimated on trade data. The value matters for the
identification of $p$ rather than for the transmission mechanism: at $\eta<1$
external balance identifies the terms of trade only weakly, which is one of the
reasons the steady state is imposed symmetrically.

\paragraph{Financial intermediation.}
The block contributes the exit/payout share $f_X$, the divertability parameter
$\lambda_X$, the entrant transfer $\omega_X$ and the bankers' discount factor
$\beta^{\rm int}_X$, alongside the deposit rate $\bar r$ it shares with the
household problem. Only the first and last are set directly: $f_X=0.04$, an
expected banker horizon of $25$ quarters, and
$\beta^{\rm int}_X=1/(1+\bar r)=0.9970$. The latter is the level term of the
banker's kernel \eqref{eq:branch-sdf}, whose state contingency --- and hence the
risk premium in \eqref{eq:bond-decomposition} --- comes entirely from the
branch-contingent GHH composite; it is not the household discount factor, which
satisfies $\beta_X(1+\bar r)<1$. Since $\bar\varphi$ falls below unity once
$\beta^{\rm int}\ll1/(1+\bar r)$, muting the franchise channel, it has to move
with $\bar r$ rather than be fixed independently.

The pair $(\lambda_X,\omega_X)$ is not assigned but solved, so that the steady
state reproduces two targets: leverage $\bar\theta$ and the intermediation wedge
$\bar\varsigma\equiv\bar r^{k}-\bar r=\lambda_X\bar\mu/\bar\Omega$. With
\eqref{eq:IC} binding, Proposition~\ref{prop:mu} and the net-worth recursion
\eqref{eq:nw-lom} collapse onto
\begin{align}
\bar\varphi&=\frac{\bar\Omega(1+\bar r)}{1-\bar\mu},
&
\bar\Omega&=\beta^{\rm int}\bigl[f+(1-f)\bar\varphi\bigr],
&
\bar\mu&=\frac{\bar\Omega\,\bar\varsigma\,\bar\theta}{\bar\varphi} ,
\label{eq:ss-system}
\end{align}
which is triangular: $\bar\varphi$ and $\bar\Omega$ cancel from the third
equation, so the targets pin the multiplier and the franchise value is linear in
it. Imposing $\beta^{\rm int}(1+\bar r)=1$,
\begin{align}
\bar\mu&=\frac{\bar\varsigma\bar\theta}{1+\bar r+\bar\varsigma\bar\theta},
&
\bar\varphi&=\frac{f}{f-\bar\mu},
&
\lambda_X&=\frac{\bar\varphi}{\bar\theta},
&
\omega_X&=\frac{\mathcal{D}}{\bar\theta}-(1-f)\,\bar\varsigma,
\label{eq:ss-closedform}
\end{align}
with $\mathcal{D}\equiv1-(1-f)(1+\bar r)$ the discount in the net-worth
accumulation identity. Three implications are worth recording. The multiplier
depends only on the product $\bar\varsigma\bar\theta$, the excess return on
intermediary equity, so leverage and the wedge are not separately identified by
it. Stationary net worth requires $\mathcal{D}>0$, i.e.\ $f>\bar r/(1+\bar r)$.
And $\omega_X>0$ bounds the admissible target space,
$\bar\varsigma\bar\theta<\mathcal{D}/(1-f)$, here $3.87\%$ per quarter against a
calibrated $0.10\%$.

The single-$\lambda$ assumption \citep[eq.~3]{bocola2016pass} gives a check on
the solved steady state. Every class earns the same excess return
$\lambda_X\bar\mu/\bar\Omega=\bar\varsigma$, so the net worth implied by the
binding constraint, $\lambda_X\bar{\mathcal A}/\bar\varphi$, and the net worth
implied by accumulation,
$\mathcal{D}^{-1}\sum_j[(1-f)(\bar R_j-\bar R)+\omega_X]\bar a_j$, are both
proportional to the divertable base. Their equality is then a scalar equation in
$\bar r^{k}$, independent of the scale \emph{and} the composition of the balance
sheet, and solving it returns $\bar r^{k}=\bar r+\bar\varsigma$ and
$\bar\theta_{ss}=\bar\theta$ to machine precision for any portfolio.

\paragraph{The two targets.}
Leverage is set to $\bar\theta=5$ and, as stated, $f=0.04$, both posterior means
of \citet{bocola2016pass}. Neither has a clean accounting counterpart: the
intermediary here holds the entire capital stock, the domestically-held
sovereign, a cross-border sovereign position and the working-capital book, so
$\bar\theta$ is leverage on that consolidated portfolio rather than the book
assets-to-equity ratio of the observed banking system, and $f$ is a payout rate,
not an exit frequency.

The wedge $\bar\varsigma$ is the least identified parameter of the model and we
treat it as such. In the deterministic steady state used for calibration the
default probability is off, so $\bar\varsigma$ is the \emph{entire} spread of the
sovereign yield over the risk-free rate, and any measured spread is an upper
bound on it. The Greek--German ten-year differential averaged $50$ basis points
over $1999\text{Q}1$--$2007\text{Q}4$.\footnote{FRED series
\texttt{IRLTLT01GRM156N} and \texttt{IRLTLT01DEM156N}; the $2010\text{Q}1$--$2012\text{Q}2$
average is $1{,}248$ basis points, peaking at $2{,}924$ in February 2012.} We set
$\bar\varsigma=8$ basis points per annum following \citet{bocola2016pass},
implying an excess return on intermediary equity
$\bar\varsigma\bar\theta=40$ basis points per annum and a barely-binding
constraint, $\bar\mu=0.0010$.

Two qualifications. First, the calibration target and the model's average spread
are different objects: at the stochastic rest point $\pi^{d}(\bar s)=0.1\%$ per
quarter is priced, which at $\varrho_D=0.45$ adds some $22$ basis points per
annum of default compensation, taking the rest-point spread to roughly $30$ basis
points before the covariance premium of \eqref{eq:bond-decomposition}. It is that
object, not $\bar\varsigma$, that the pre-crisis differential should be compared
to. Second, $\bar\varsigma\bar\theta$ is the rate at which intermediaries rebuild
equity out of retained earnings and therefore governs the persistence of any
balance-sheet shock: at $\bar\varsigma=200$ basis points the excess return on
equity reaches $10\%$ per annum, the equity loss generated by our risk experiment
is recouped within about two quarters, and the output response turns positive
from the third year. The small wedge is what makes transmission run through the
level of net worth and the price of the sovereign claim. Its corollary is that at
$\bar\mu=0.0010$ the steady state sits near the boundary of the
occasionally-binding region and the constraint goes slack some four quarters into
the experiment, so the sustained contraction is carried by $q^{D}$ and $n_D$
rather than by a persistently elevated credit spread. We report sensitivity to
$\bar\varsigma$ in Section~[\textsc{tbc}].

\paragraph{Interest rates.}
The risk-free rate is the euro-area real short rate: deflating the three-month
interbank rate by headline HICP inflation over
$1999\text{Q}1$--$2007\text{Q}4$ gives $1.11\%$ per annum, and we set
$\bar r=0.003$ per quarter ($1.20\%$ p.a.).\footnote{FRED series
\texttt{IR3TIB01DEM156N} and \texttt{CP0000EZ19M086NEST}. The pre-crisis window
is the relevant one: the $2010\text{Q}1$--$2012\text{Q}2$ average real rate is
$-1.19\%$ p.a. and does not represent the stationary rate around which the crisis
experiment is run.} Deposits are own-good claims at national rates tied by
$(1+r_{D,t})=(1+r_{F,t})p_{t+1}/p_t$, the flexible-price image of one union
policy rate with national inflation differentials; the cross-border deposit
position this permits is the margin along which national saving and investment
decouple, and it is zero at the symmetric steady state.

With a single $\lambda_X$ every class must deliver the same excess return, so the
steady-state cross-section of returns collapses onto $\bar r$ and
$\bar\varsigma$:
\begin{equation}
\bar r^{k}=\bar r^{b}_{D}=\bar r^{b}_{F}=\bar r^{wc}=\bar r+\bar\varsigma
=1.28\%\ \text{p.a.},
\qquad
\bar q^{X}=\frac{\delta_b}{\bar r+\delta_b+\bar\varsigma}=0.9459 ,
\label{eq:rate-structure}
\end{equation}
with $\bar r^{n}=\bar r+\bar\theta\bar\varsigma=1.60\%$ per annum on
intermediary equity. The model thus matches one average spread and generates no
term premium, which is why the observed sovereign differential is not a clean
counterpart to $\bar\varsigma$.

\paragraph{The sovereign claim.}
Both governments issue \citet{Hatchondo2009} perpetuities whose stock decays at
rate $1-\delta_b$, so that Macaulay duration is
$(1+\bar r^{b})/(\bar r^{b}+\delta_b)$ quarters, evaluated at the bond's own
yield. Setting $\delta_b=0.056$ delivers a duration of $16.9$ quarters, or $4.2$
years, matching the duration of Greek marketable central-government debt on the
eve of the crisis [\textsc{tbc}].\footnote{Public Debt Management Agency,
\emph{Public Debt Bulletin}; duration, not average residual maturity, is the
correct counterpart: it is the elasticity of the bond price to the discount rate,
and it is that elasticity which governs the mark-to-market loss intermediaries
take when risk is priced. The long duration is load-bearing: at $\delta_b=0.25$
the repricing of the stock shrinks by roughly a factor of six and the
pass-through mechanism largely disappears.}

The \emph{size} of the sovereign position is disciplined by intermediary exposure
rather than by the public debt-to-GDP ratio: only the bank-intermediated portion
of the debt appears in the model. We set the
outstanding stock $\bar B_X=0.98$ --- $24.5\%$ of annual output at face value,
$23.2\%$ at market value --- which puts domestic sovereign holdings at $7.1\%$ of
the domestic intermediary's assets and total sovereign holdings, including the
cross-border leg, at $8.9\%$ [\textsc{tbc}].\footnote{The exposure share is the
figure reported by \citet{bocola2016pass}; replacing it with the corresponding
Greek supervisory disclosure (European Banking Authority, 2011 EU-wide stress test
and capital exercise) is a first-order change, because the exposure share scales
the entire balance-sheet channel linearly.} Each country's intermediary holds
$80\%$ of its own sovereign and $20\%$ of the other's, the cross-border leg along
which the shock travels. The distinction is decisive: reading the exposure ratio
against GDP rather than against intermediary assets overstates the position
roughly threefold, enough for the fiscal relief from a default to outweigh
intermediaries' losses and reverse the sign of the mechanism. The rest of the
balance sheet is capital ($85.5\%$ of assets) and the working-capital book
($5.6\%$).

Government purchases are set to $G_X=0$, so the fiscal block finances only the
coupon net of new issuance; steady-state taxes are correspondingly small at
$0.30\%$ of quarterly output. The Bohn rule \eqref{eq:bohn} is imposed as a
constant elasticity of the tax level with respect to the surviving debt stock,
with $\gamma_\tau=1$, the value in \citet{bocola2016pass}. The elasticity form
matters for the default event: written instead as a response to the debt
\emph{level}, a $55\%$ haircut would hand households a tax windfall worth tens of
percent of GDP and again invert the sign of the mechanism.

\paragraph{The default event.}
The recovery rate in the feared event is taken from the realised restructuring.
The March 2012 PSI exchange imposed a face-value reduction of $53.5\%$, and
\citet{zettelmeyer2013} estimate the net-present-value haircut at $59$--$65\%$
depending on the discount rate. We set the recovery rate to $\varrho_D=0.45$, at
the conservative end of that range, applied to the coupon and to the continuation
value alike as in \eqref{eq:bond-payoff}. There is no exogenous output cost of
default, no capital-quality loss and no recapitalisation rule: the default-state
recession comes entirely from \eqref{eq:ng}--\eqref{eq:mu-closed} operating on a
smaller asset base. The severity of the default state is thus an outcome of the
intermediary block, and the risk premium cannot be tuned through it.

\paragraph{The sovereign risk process.}
The latent factor $s_t$ in \eqref{eq:s-process} is calibrated so that the
rest-point default probability is $\pi^{d}(\bar s)=0.1\%$ per quarter, i.e.\
$\bar s=\log(0.001/0.999)=-6.91$, with persistence $\rho_s=0.95$ and innovation
standard deviation $\sigma_s=0.63$, both posterior means from
\citet{bocola2016pass}; the unconditional standard deviation of $s$ is
$\sigma_s/\sqrt{1-\rho_s^{2}}=2.02$. The process is calibrated to a level and a
persistence, not to a one-quarter jump: at $\rho_s=0.95$ crisis-level default
probabilities are reached by wandering, and asking a single innovation to lift
$\pi^{d}$ from $0.1\%$ to $2\%$ would require $\sigma_s\approx1.5$, outside both
the estimated process and any workable collocation box.

The baseline experiment sets the priced probability $\pi^{d}_t>0$ with the
realised indicator $d_t\equiv0$ throughout: risk is priced but never realised,
and all dynamics are generated by the anticipation of an event that does not
occur. The impulse we report raises $\pi^{d}$ on impact from $0.10\%$ to $1.98\%$
per quarter --- a shift of $3.0$ in $s$, or about $1.5$ unconditional standard
deviations --- after which it decays at $\rho_s$. For comparison we also report a
productivity experiment, a $1\%$ innovation to $Z_D$ decaying at $\rho_z=0.9$.

\begin{table}[t]
\centering
\caption{Implied steady-state objects}
\label{tab:ss-implied}
\begin{tabular}{lcl}
\hline\hline
Object & Value & Determined by \\
\hline
Divertable share $\lambda_X$              & $0.2051$ & solved, \eqref{eq:ss-closedform} \\
Entrant transfer $\omega_X$               & $0.0072$ & solved, \eqref{eq:ss-closedform} \\
Franchise value $\bar\varphi$             & $1.0254$ & solved, \eqref{eq:ss-closedform} \\
IC multiplier $\bar\mu$                   & $0.0010$ & solved, \eqref{eq:ss-closedform} \\
Household discount factor $\beta_X$       & $0.9945$ & deposit clearing, \eqref{eq:beta-clearing} \\
Bond price $\bar q^{X}$                   & $0.9459$ & \eqref{eq:rate-structure} \\
Return on equity $\bar r^{n}$ (p.a.)      & $1.60\%$ & $\bar r+\bar\theta\bar\varsigma$ \\
Capital--output ratio (annual)            & $2.22$   & firm FOC at $\bar r^{k}$ \\
Investment share $\delta\bar K/\bar Y$    & $22.2\%$ & firm FOC at $\bar r^{k}$ \\
Deposits / annual output                  & $2.07$   & $(\bar\theta-1)\bar n$ \\
Sovereign share of intermediary assets    & $7.1\%$  & calibration target on $\bar B$ \\
\hline\hline
\end{tabular}

\smallskip
\begin{minipage}{0.92\textwidth}
\footnotesize \emph{Note.} None of these objects is assigned; each is implied by
the parameters of Section~\ref{sec:calibration} and reported at the symmetric
steady state, in which $\bar p=1$, $\bar N_X=1$ and $\bar Y_X=1$.
\end{minipage}
\end{table}

\subsection{Solution Method}
\label{sec:solution}

The model is solved globally, as recursive decision rules: neither the kink in
\eqref{eq:mu-closed} nor the covariance premium in
\eqref{eq:bond-decomposition} survives a local approximation.

The aggregate state is the seven-dimensional vector
$\bigl[K_D,K_F,P_D,P_F,B_D,s,Z_D\bigr]$: the two capital stocks, the two
intermediaries' gross obligation states \eqref{eq:P-state}, the risky debt stock,
the sovereign-risk factor and the productivity state. Decision rules are
approximated by Chebyshev polynomials on a Smolyak sparse grid over a box centred
on the steady state, with half-widths of $3\%$ on the capital stocks and on
$Z_D$, $25\%$ on the obligation states, $30\%$ on the debt stock, and $\pm4.35$
on $s$ --- the last being $\pm2.16$ unconditional standard deviations of the risk
process, matching the coverage in \citet{bocola2016pass}. The baseline risk
experiment uses an isotropic Smolyak level $\mu=1$, i.e.\ $15$ collocation
points; the productivity experiment, which has no risk dimension to resolve, uses
$\mu=2$ ($113$ points).

At each collocation point the seven market-clearing unknowns
$[N_D,N_F,K_D',K_F',r_D,r_F,p]$ are solved given the previous iterate's rules as
the continuation, with the multiplier taken from its closed form
\eqref{eq:mu-closed}; the solver then iterates on the rules to convergence. The
conditional expectation \eqref{eq:expectation} is evaluated as a genuine
multi-branch quadrature: a seven-node Gauss--Hermite rule over the innovation to
$s$, crossed with the default fork, with the default branch evaluated on the same
fitted decision rules at a reachable next-period state rather than on a frozen
stand-in economy. Setting $\pi^{d}\equiv0$ recovers the risk-neutral model
exactly, which we use as a regression test. Solution accuracy is reported as
Euler errors along a long simulation of the ergodic set. [\textsc{tbc}: report the
error statistics and the box-escape frequency here.]
